% ───────────────────────────────────────────────────────────
% Bilingual LaTeX Template for academic-translation
%
% 用法：
%   xelatex bilingual.tex   # 必须用 xelatex 以支持 xeCJK
%
% 字体：
%   英文：Latin Modern (LaTeX 默认)
%   中文：STSong（macOS 自带） / SimSun（Windows） / Noto Serif CJK SC（Linux）
%
% 契约：bilingual_tex 输出契约 + xeCJK 中文支持
% ───────────────────────────────────────────────────────────

\documentclass[11pt,a4paper]{article}

% 中文支持（xeCJK）
\usepackage{xeCJK}
% 多平台字体 fallback：依次尝试
\IfFontExistsTF{STSong}{
  \setCJKmainfont{STSong}
}{
  \IfFontExistsTF{SimSun}{
    \setCJKmainfont{SimSun}
  }{
    \setCJKmainfont{Noto Serif CJK SC}
  }
}

% 双栏并行（paracol）
\usepackage{paracol}
\setlength{\columnsep}{1.5em}
\setlength{\columnseprule}{0.4pt}

% 数学
\usepackage{amsmath, amssymb, amsthm}

% 引用 / 链接
\usepackage[backend=bibtex]{biblatex}
\usepackage{hyperref}
\hypersetup{colorlinks=true, linkcolor=blue, citecolor=blue, urlcolor=blue}

% 页面
\usepackage[margin=2cm]{geometry}

% 图 / 表
\usepackage{graphicx}
\usepackage{booktabs}

% 元数据
\title{{Paper Title (English)} \\ {论文标题（中文）}}
\author{Translated by academic-translation Skill}
\date{\today}

% Bibliography file（用户提供）
% \addbibresource{references.bib}

\begin{document}
\maketitle

% ───────────────────── Abstract / 摘要 ─────────────────────
\begin{paracol}{2}
\section*{Abstract}
We propose a novel method that achieves state-of-the-art on
ImageNet \cite{deng2009imagenet}, with $\mathbf{x}_t$ defined as
\begin{equation}
  \mathbf{x}_t = f(\mathbf{x}_{t-1}; \theta).
\end{equation}
Extensive experiments demonstrate ...

\switchcolumn

\section*{摘要}
我们提出一种新颖的方法，在 ImageNet \cite{deng2009imagenet}
上取得最先进的表现，其中 $\mathbf{x}_t$ 定义为
\begin{equation}
  \mathbf{x}_t = f(\mathbf{x}_{t-1}; \theta).
\end{equation}
大量实验表明……
\end{paracol}

% ───────────────────── §1 Introduction / 引言 ─────────────────────
\begin{paracol}{2}
\section{Introduction}
Deep learning has achieved remarkable success
\cite{lecun2015deep}, including image recognition,
machine translation, and natural language understanding.

% 段落 §1-p2
However, despite these advances, several challenges remain:
\begin{itemize}
  \item Challenge A: ...
  \item Challenge B: ...
\end{itemize}

\switchcolumn

\section{引言}
深度学习已在图像识别、机器翻译和自然语言理解等多个领域取得显著成功
\cite{lecun2015deep}。

% 段落 §1-p2
然而，尽管这些进展可观，仍存在若干挑战：
\begin{itemize}
  \item 挑战 A：……
  \item 挑战 B：……
\end{itemize}
\end{paracol}

% ───────────────────── §2 Related Work / 相关工作 ─────────────────────
% 按相同 paracol 模式继续

% ───────────────────── References / 参考文献 ─────────────────────
% \printbibliography

\end{document}

% ───────────────────────────────────────────────────────────
% 模板使用说明：
%
% 1. 替换 {Paper Title} 为实际标题
% 2. 用 academic-translation 输出的 English 段落填左栏
% 3. 用 step3_polished 输出的 Chinese 段落填右栏
% 4. **MUST**：所有 \cite{}、\ref{}、$...$、\eqref{} 在左右栏一字不差
% 5. 编译：xelatex bilingual.tex
%
% 字体不存在时的降级：
% - 如果三个 CJK 字体都不存在，编译会失败。可改用：
%   \setCJKmainfont{Songti SC}  (macOS)
%   \setCJKmainfont{Noto Sans CJK SC}  (任何平台)
%
% Docker 编译（无本地 xelatex 时）：
% docker run --rm -v $PWD:/work -w /work \
%   ghcr.1ms.run/xu-cheng/texlive-debian:20260101 \
%   xelatex bilingual.tex
% ───────────────────────────────────────────────────────────
